\chapter{Paysage et pipelines de donn\'ees}
\label{ch:data-landscape}

\begin{quote}
\emph{<< La donn\'ee est le nouveau p\'etrole --- mais comme le p\'etrole, elle doit \^etre raffin\'ee avant d'\^etre utile. >>}
--- Clive Humby
\end{quote}

% --- hook ---
Vous arrivez votre premier jour comme data scientist dans une banque de Cotonou. On vous tend une clé USB : << voilà cinq ans de données de transactions, fais-nous un modèle de détection de fraude pour vendredi. >> Vous ouvrez le CSV. Trois colonnes ont des en-têtes en wolof. Une autre contient à la fois des dates au format français et au format ISO. Le champ << montant >> a des valeurs négatives qu'aucun virement réel n'expliquerait. C'est à ce moment précis que vous comprenez ce que veulent dire les enquêtes sectorielles quand elles annoncent que les data scientists passent 60 à 80\% de leur temps à nettoyer. Ce cours traite cette phase comme ce qu'elle est : le cœur du métier, pas son préambule.
\section{Qu'est-ce que le pr\'etraitement des donn\'ees ?}

Le pr\'etraitement des donn\'ees est l'ensemble des transformations appliqu\'ees aux donn\'ees brutes avant analyse ou mod\'elisation. Il fait le pont entre les donn\'ees que vous \emph{avez} et celles dont vous \emph{avez besoin}.

Un pipeline typique de pr\'etraitement inclut :
\begin{enumerate}
  \item \textbf{Chargement} --- lire les donn\'ees depuis fichiers, bases de donn\'ees ou API.
  \item \textbf{Inspection} --- comprendre structure, types et qualit\'e.
  \item \textbf{Nettoyage} --- g\'erer valeurs manquantes, doublons et erreurs.
  \item \textbf{Transformation} --- encoder, mettre \`a l'\'echelle, cr\'eer des variables.
  \item \textbf{Validation} --- confirmer que la sortie r\'epond aux besoins en aval.
\end{enumerate}

\begin{datatip}
Selon les enqu\^etes sectorielles, les praticiens consacrent 60--80\% du temps de projet au pr\'etraitement. Ma\^itriser cette phase est la comp\'etence la plus levier de la science des donn\'ees.
\end{datatip}

% ============================================================
\section{Types et structures de donn\'ees}

\subsection{Donn\'ees structur\'ees, semi-structur\'ees et non structur\'ees}

\begin{itemize}
  \item \textbf{Structur\'ees :} tables avec lignes et colonnes (CSV, bases SQL, Excel).
  \item \textbf{Semi-structur\'ees :} donn\'ees \`a sch\'emas souples (JSON, XML, YAML).
  \item \textbf{Non structur\'ees :} texte libre, images, audio, vid\'eo.
\end{itemize}

Ce cours se concentre sur les donn\'ees structur\'ees et semi-structur\'ees, avec un chapitre (Chapitre~7) consacr\'e au texte.

\subsection{Types de variables}

\begin{minted}{python}
import pandas as pd

# Charger un dataset r\'eel pour inspecter les types
url = ("https://raw.githubusercontent.com/datasciencedojo/"
       "datasets/master/titanic.csv")
df = pd.read_csv(url)
print(df.dtypes)
\end{minted}

Sortie (abr\'eg\'ee) :
\begin{minted}{text}
PassengerId      int64
Survived         int64
Pclass           int64
Name            object
Sex             object
Age            float64
SibSp            int64
Fare           float64
Embarked        object
\end{minted}

\begin{preprocesstip}
Le dtype \texttt{object} dans pandas indique g\'en\'eralement des cha\^ines, mais peut aussi cacher des types m\'elang\'es. V\'erifiez toujours avec \texttt{df[col].apply(type).value\_counts()}.
\end{preprocesstip}

\subsection{Num\'erique vs cat\'egoriel}

\begin{minted}{python}
numerical = df.select_dtypes(include="number").columns.tolist()
categorical = df.select_dtypes(include="object").columns.tolist()

print(f"Numerical columns ({len(numerical)}): {numerical}")
print(f"Categorical columns ({len(categorical)}): {categorical}")
\end{minted}

% ============================================================
\section{Formats de donn\'ees courants}

\begin{longtable}{lllp{5cm}}
\toprule
\textbf{Format} & \textbf{Extension} & \textbf{Lecteur pandas} & \textbf{Remarques} \\
\midrule
CSV & \texttt{.csv} & \texttt{read\_csv} & Le plus courant ; attention au d\'elimiteur et \`a l'encodage \\
Excel & \texttt{.xlsx} & \texttt{read\_excel} & N\'ecessite \texttt{openpyxl} ; support multi-feuilles \\
JSON & \texttt{.json} & \texttt{read\_json} & Les structures imbriqu\'ees demandent \texttt{json\_normalize} \\
Parquet & \texttt{.parquet} & \texttt{read\_parquet} & Colonnaire, rapide, pr\'eserve les types \\
SQL & --- & \texttt{read\_sql} & N\'ecessite une connexion \texttt{sqlalchemy} \\
\bottomrule
\end{longtable}

% ============================================================
\section{Premier regard sur un dataset r\'eel}

Nous utiliserons le dataset Melbourne Housing dans plusieurs chapitres :

\begin{minted}{python}
# Dataset Melbourne Housing
melb = pd.read_csv("melb_data.csv")

# Shape : lignes x colonnes
print(f"Shape: {melb.shape}")

# Trois premi\`eres lignes
print(melb.head(3))

# Statistiques de r\'esum\'e
print(melb.describe())

# Valeurs manquantes par colonne
print(melb.isnull().sum().sort_values(ascending=False).head(10))
\end{minted}

\begin{minted}{python}
# Rapport rapide de qualit\'e
def data_quality_report(df):
    """Generate a quick quality report for any DataFrame."""
    report = pd.DataFrame({
        "dtype": df.dtypes,
        "non_null": df.notnull().sum(),
        "null_count": df.isnull().sum(),
        "null_pct": (df.isnull().sum() / len(df) * 100).round(1),
        "n_unique": df.nunique(),
        "sample_value": df.iloc[0]
    })
    return report.sort_values("null_pct", ascending=False)

print(data_quality_report(melb))
\end{minted}

\begin{datatip}
Faites toujours un rapport de qualit\'e \emph{avant} toute analyse. Les cinq nombres dont vous avez imm\'ediatement besoin : shape, dtypes, comptes de nuls, comptes d'unicit\'e, et statistiques de r\'esum\'e.
\end{datatip}

% ============================================================
\section{Le workflow de pr\'etraitement}

\begin{minted}{python}
# Workflow minimal de bout en bout
import pandas as pd
import numpy as np

# 1. Charger
df = pd.read_csv("melb_data.csv")

# 2. Inspecter
print(df.info())

# 3. Nettoyer : doublons, valeurs manquantes
df = df.drop_duplicates()
df["Car"] = df["Car"].fillna(df["Car"].median())
df = df.dropna(subset=["Price"])

# 4. Transformer : encoder cat\'egoriel, mettre \`a l'\'echelle num\'erique
df["Type_encoded"] = df["Type"].map({"h": 0, "u": 1, "t": 2})

# 5. Valider
assert df.isnull().sum().sum() == 0 or True  # relax\'e pour la d\'emo
print(f"Clean shape: {df.shape}")
\end{minted}

% ============================================================
\section{R\'evision pandas : op\'erations essentielles}

\subsection{S\'election et filtrage}

\begin{minted}{python}
# S\'electionner des colonnes
prices = melb[["Suburb", "Price", "Rooms"]]

# Filtrer des lignes
expensive = melb[melb["Price"] > 2_000_000]
south = melb[melb["Regionname"] == "Southern Metropolitan"]

# Combin\'e
big_south = melb[(melb["Regionname"] == "Southern Metropolitan") &
                 (melb["Rooms"] >= 4)]
print(f"Large southern houses: {len(big_south)}")
\end{minted}

\subsection{Groupby et agr\'egation}

\begin{minted}{python}
# Prix moyen par r\'egion
region_stats = (melb.groupby("Regionname")["Price"]
                .agg(["mean", "median", "count"])
                .sort_values("median", ascending=False))
print(region_stats)
\end{minted}

\begin{warningbox}
Pandas \texttt{groupby} \'elimine silencieusement les cl\'es \texttt{NaN} par d\'efaut. Si votre colonne de groupement a des valeurs manquantes, elles seront exclues du r\'esultat. Utilisez \texttt{dropna=False} pour les inclure.
\end{warningbox}

% ============================================================
\section{Aper\c{c}u des outils et biblioth\`eques}

\begin{minted}{python}
# Pile de pr\'etraitement de base
import pandas as pd               # manipulation de donn\'ees
import numpy as np                 # op\'erations num\'eriques
import matplotlib.pyplot as plt    # visualisation
import seaborn as sns              # graphiques statistiques
import missingno as msno           # visualisation des manquants
from sklearn.preprocessing import StandardScaler, OneHotEncoder
from sklearn.impute import SimpleImputer, KNNImputer
from sklearn.pipeline import Pipeline
from sklearn.compose import ColumnTransformer
\end{minted}

% ============================================================
\section{Exercices}

\begin{exercise}
\begin{enumerate}
  \item T\'el\'echargez le dataset Melbourne Housing. \'Ecrivez un rapport de qualit\'e montrant : nombre de lignes, colonnes, dtype de chaque colonne, pourcentage de valeurs manquantes, et nombre de valeurs uniques.
  \item Chargez le dataset Gapminder (\texttt{gapminder.csv} ou via \texttt{plotly.express}). Identifiez les colonnes num\'eriques, cat\'egorielles et celles \`a valeurs manquantes.
  \item \'Ecrivez une fonction \texttt{detect\_mixed\_types(df)} qui v\'erifie chaque colonne pour des types Python m\'elang\'es et renvoie la liste des colonnes probl\'ematiques.
  \item Comparez le chargement du dataset Titanic en CSV vs Parquet. Mesurez taille de fichier et temps de chargement avec \texttt{\%timeit}.
\end{enumerate}
\end{exercise}

% ============================================================
\section{R\'esum\'e du chapitre}

\begin{itemize}
  \item Le pr\'etraitement transforme les donn\'ees brutes en donn\'ees pr\^etes pour l'analyse et consomme la majorit\'e du temps de projet.
  \item Les donn\'ees existent en formats structur\'es, semi-structur\'es et non structur\'es ; pandas g\`ere les deux premiers.
  \item Tout workflow de pr\'etraitement suit cinq \'etapes : charger, inspecter, nettoyer, transformer, valider.
  \item Un rapport de qualit\'e rapide (shape, dtypes, nuls, unicit\'es, statistiques) devrait \^etre votre premi\`ere action sur n'importe quel nouveau dataset.
  \item La pile Python de base pour le pr\'etraitement : pandas, numpy, scikit-learn, matplotlib, missingno.
\end{itemize}
